\subsection{Etapa 2}
\begin{frame}{Etapa 2}
    A próxima estapa estima a fase $\theta$
    \begin{equation}
        \theta_{pred} = NN_{\theta} (Y_{obs};\phi_{\theta})
    \end{equation}
    A arquitetura possui:
    \begin{itemize}
        \item 4 camadas convolucionais;
        \item 2 fully connected layers.
    \end{itemize}
    A primeira convolução usa kernel \(29\times1\) e stride 14.

    As seguintes usam \(4\times1\) e stride 2.

    Depois ocorre flattening e duas camadas fully connected, resultando em um único valor de fase.

    O método usa a transformada de Hilbert para girar $w_0$ no domínio de fase por $\theta$
    \begin{equation}
        w_{pred} = Re[(w_0 + j Hilbert(w_0)) e^{-j\theta_{pred}}]
    \end{equation}
\end{frame}